%! Exponential and Logarithmic Equations and Inequalities — Unit 2, Lesson 2.13 — Homework
\documentclass[10pt]{article}
\usepackage{precalculus-article}
\usepackage{precalculus-boxes}
\newcommand{\TallMath}[1]{$\displaystyle #1\rule[-1.4em]{0pt}{3.2em}$}

\begin{document}

\pageheader{Unit 2, Lesson 2.13}{Homework}
\namedateperiod

\vspace{0.1in}

\begin{enumerate}[leftmargin=1.8em, itemsep=14pt]

\item Solve each equation for $x$. Where necessary, give both exact and decimal forms
      (rounded to 3 decimal places).

\begin{enumerate}[leftmargin=1.5em, itemsep=8pt, label=\alph*.]
  \item $3^x = 243$

  \vspace{4pt}
  $x = $ \blank{4cm}

  \item $2^x = 50$

  \vspace{4pt}
  Exact: $x = $ \blank{4cm} \qquad Decimal: $x \approx $ \blank{2.5cm}

  \item $7^x = \dfrac{1}{49}$

  \vspace{4pt}
  $x = $ \blank{4cm}
\end{enumerate}

\item Solve each logarithmic equation. Show all steps and check for extraneous solutions.

\begin{enumerate}[leftmargin=1.5em, itemsep=10pt, label=\alph*.]
  \item $\log_2(x+3) = 5$

  \vspace{4pt}
  Work: \writeline

  \vspace{4pt}
  Solution: $x = $ \blank{2.5cm} \qquad Valid? \blank{2cm}

  \item $\log_3 x + \log_3(x-8) = 2$

  \vspace{4pt}
  Step 1 (combine): \writeline

  \vspace{4pt}
  Step 2 (exponential form): \writeline

  \vspace{4pt}
  Step 3 (solve quadratic): \writeline

  \vspace{4pt}
  Possible solutions: $x = $ \blank{2cm} or $x = $ \blank{2cm}

  \vspace{4pt}
  Check for extraneous solutions: \writeline

  \vspace{4pt}
  Final answer: $x = $ \blank{2.5cm}
\end{enumerate}

\item A population model is $P(t) = 1200 \cdot (1.04)^t$, where $t$ is in years.
      For what value of $t$ does $P(t) = 2000$? Show work using logarithms.

\vspace{4pt}
Step 1: \writeline

\vspace{4pt}
Step 2: \writeline

\vspace{4pt}
Step 3: \writeline

\vspace{4pt}
$t = $ \blank{5cm} (exact) $\approx $ \blank{2.5cm} years

\item Rewrite each expression using the natural base identity $b^x = e^{(\ln b)\cdot x}$.

\begin{enumerate}[leftmargin=1.5em, itemsep=6pt, label=\alph*.]
  \item $5^x = e^{(\;\blank{3cm}\;)\cdot x}$

  \item $(0.7)^t = e^{(\;\blank{3cm}\;)\cdot t}$
\end{enumerate}

\item Find the inverse of $f(x) = 2 \cdot 5^{x+3} - 1$. Show each step. Then state the
      domain and range of both $f$ and $f^{-1}$.

\vspace{6pt}
\writelines{4}

\vspace{4pt}
$f^{-1}(x) = $ \blank{6cm}

\vspace{6pt}
Domain of $f$: \blank{3.5cm} \qquad Range of $f$: \blank{3.5cm}

\vspace{4pt}
Domain of $f^{-1}$: \blank{3.5cm} \qquad Range of $f^{-1}$: \blank{3.5cm}

\item \textbf{AP-style multiple choice.} Which value of $x$ satisfies
      $\log_2(x+4) + \log_2 x = 5$?

\vspace{6pt}
\begin{enumerate}[leftmargin=2em, itemsep=4pt, label=(\Alph*)]
  \item $x = -8$
  \item $x = 4$
  \item $x = 2$
  \item $x = 8$
\end{enumerate}

\end{enumerate}

\vspace{0.1in}

\begin{extensionbox}
\textbf{Extension:} Solve the inequality $\log_3(x+1) \geq 2$. Express the solution as an
inequality and justify why $x = -1$ is not in the solution set.

\writelines{3}
\end{extensionbox}

\vspace{0.1in}

\begin{reflectionbox}
What is one technique from today's lesson you feel confident about, and one situation
(equation type) where you still want more practice?

\writelines{2}
\end{reflectionbox}

\end{document}
